\documentclass[12pt,letterpaper]{article}
\usepackage[utf8]{inputenc}
\usepackage[spanish]{babel}
\usepackage{amsmath}
\usepackage{amsfonts}
\usepackage{amssymb}
\usepackage[margin=2.5cm]{geometry}

\begin{document}

\begin{center}
    \Large \textbf{Evaluación de Examen 2: Unidad II} \\
    \large Física para Ciencias de la Salud (005-1713) \\
    \vspace{0.5cm}
    \textbf{Estudiante:} Daniela Rivero \quad \textbf{C.I.:} 31.435.612
    \vspace{0.3cm} \\
    \large \textbf{Calificación Final: 10/10 puntos}
\end{center}

\vspace{0.5cm}

\section*{Análisis de Resultados}

\subsection*{1. Cinemática (Aceleración media)}
\textbf{Procedimiento y resultado:} Estructura excelentemente los datos ($v_0 = 0 \text{ m/s}$, $v_f = 0,6 \text{ m/s}$, $t = 0,15 \text{ s}$) y plantea la fórmula del cambio de velocidad respecto al tiempo. El desarrollo de $0,6 / 0,15$ es correcto, obteniendo el resultado de $4 \text{ m/s}^2$ con el análisis dimensional exacto. \\
\textbf{Veredicto:} \textbf{Correcto} (2/2 pts).

\subsection*{2. Dinámica (Leyes de Newton)}
\textbf{Procedimiento y resultado:} Extrae la fuerza y la masa. Plantea la fórmula de la Segunda Ley de Newton con un despeje correcto de la aceleración ($a = \frac{F}{m}$). La división de $150 \text{ N}$ entre $25 \text{ kg}$ arroja un resultado perfecto de $6 \text{ m/s}^2$. \\
\textbf{Veredicto:} \textbf{Correcto} (2/2 pts).

\subsection*{3. Trabajo Mecánico}
\textbf{Procedimiento y resultado:} Identifica y plasma la fórmula base del trabajo mecánico ($W = F \cdot d$). La sustitución numérica de $400 \text{ N} \cdot 0,8 \text{ m}$ la conduce al cálculo correcto de $320$. Añade acertadamente las unidades correspondientes en Joules. \\
\textbf{Veredicto:} \textbf{Correcto} (2/2 pts).

\subsection*{4. Energía Potencial}
\textbf{Procedimiento y resultado:} Extrae los tres datos clave del enunciado (masa, gravedad estándar y altura) de forma muy ordenada. Multiplica los factores aplicando la fórmula de la energía potencial gravitatoria ($E_p = m \cdot g \cdot h$) hasta llegar a la cifra exacta de $17,64 \text{ J}$. \\
\textbf{Veredicto:} \textbf{Correcto} (2/2 pts).

\subsection*{5. Potencia Mecánica}
\textbf{Procedimiento y resultado:} Vincula a la perfección los Joules de trabajo efectuado por el corazón entre los $60 \text{ s}$ del enunciado usando la ecuación de potencia ($P = \frac{W}{t}$). El resultado aritmético es el esperado: $1,25 \text{ W}$. \\
\textbf{Veredicto:} \textbf{Correcto} (2/2 pts).

\vspace{0.5cm}
\section*{Comentarios Generales y Sugerencias}
¡Muchas felicidades por un examen perfecto! La presentación de las respuestas destaca por su claridad y nivel de detalle metodológico.
\begin{itemize}
    \item El formato implementado por la estudiante (listar los \textit{Datos}, establecer la \textit{Fórmula}, realizar la \textit{Sustitución} y redactar una frase conclusiva remarcando el resultado final) es el estándar oro esperado en reportes y guías de laboratorio.
    \item Las magnitudes y análisis dimensionales fueron manejados de forma estupenda, evidenciando que comprende bien la diferencia conceptual entre las distintas unidades del S.I. aplicadas a la salud.
    \item Se invita a mantener este excelente nivel de dedicación para las futuras evaluaciones.
\end{itemize}

\end{document}